\section{Analytic qualifications for history coordinates and horizon coherence}

This note qualifies the interpretation of the existing comparisons and corrects
an overly broad horizon-repair statement in the earlier Supplement. The
arguments are algebraic or distributional: they introduce no fitted forecasts,
new empirical scores, resampling output or changes to the archived results.

\subsection*{A. Method-specific cutpoints cannot distinguish constant rescalings}

Let \(W\) count completed fitting-label origins, let \(H\) be the forecast
horizon, and let \(p_m>0\) be the declared coefficient count for a fixed method
\(m\). Consider the two candidate coordinates
\[
c=\frac{W}{H+1},\qquad
b=\frac{W}{(H+1)p_m}=\frac{c}{p_m}.
\]
The pooled discovery classifier can distinguish these coordinates when methods
have different \(p_m\). The later use of a separate cutpoint for each method
has a different invariance. Within a method, positive constant scaling
preserves the ordering and equality groups of the coordinate values. If a
zero crossing between adjacent points \((x_0,y_0)\) and \((x_1,y_1)\) is
interpolated as
\[
r=x_0+(x_1-x_0)\frac{y_0}{y_0-y_1},
\]
replacing every \(x\) by \(x/p_m\) replaces \(r\) by \(r/p_m\). Crossing
directions, crossing counts and censoring statuses are unchanged. Applied to
each existing resampling draw separately, the same identity scales the
cutpoint-interval endpoints; zero and unbounded endpoints retain their
statuses.

Writing \(r_b=r_c/p_m\), the subsequent method-specific screen regions obey
\[
b<r_b/2\ \Longleftrightarrow\ c<r_c/2,
\qquad
b>2r_b\ \Longleftrightarrow\ c>2r_c.
\]
Their interval transfers back to history are also identical:
\[
(H+1)p_m[\ell_b,u_b]=(H+1)[\ell_c,u_c].
\]
Consequently these later screens and interval transfers cannot select between
the two normalizations when each method receives its own cutpoint. This is an
exact-arithmetic property of the definitions, rather than a claim about an
alternate floating-point execution. It does not undo the specified screen
outcomes, establish a complexity mechanism, or resolve the separate
comparison between larger and smaller forecasting procedures.

\subsection*{B. A pooled discovery root is not a fixed-horizon learning root}

The discovery construction pools a method's simulation cells across horizons
and process settings before grouping and sorting by its selected coordinate.
A crossing interpolated between neighboring points that differ in both \(W\)
and \(H\) is therefore a cutpoint of that pooled grid summary. By itself it
does not identify the amount of history at which a fixed-horizon learning
curve crosses zero. Correct execution of the interpolation rule does not
remove this distinction.

For a logical example, suppose relative loss is a function of \(W\) alone,
positive at shorter histories and negative at the longest history. Sorting
the observations by \(W/(H+1)\) can still place opposite-sign observations
from different horizons next to each other and produce an interpolated
crossing. Such a crossing does not isolate an effect of the horizon
normalization. This is an identification counterexample, not an estimated
alternative process. A grid without common coordinate values across its
horizons also supplies no direct comparison of distinct horizons at the same
coordinate.

The financial learning curves instead hold a horizon and evaluation record
fixed while changing \(W\), together with the corresponding historical
reference. Agreement with an interval transferred from the pooled discovery
root is a conditional consistency check between these different objects.
It does not turn the pooled root into a horizon-specific discovery root.
For background, \href{https://arxiv.org/pdf/2103.10948}{Viering and Loog,
Section 2.1} distinguish a fixed learning problem from averages across
problems; that distinction does not itself validate the present coordinate.

\subsection*{C. Matched completed-label pools do not match every information input}

The matched procedures share the specified rolling pool of completed fitting
labels and the refit clock. Conditional features additionally consult
issue-known raw lookback observations, while the regression's penalty
selection uses its specified validation segment. Thus \(W\) counts mature
response origins rather than every raw observation consulted for predictors
or every observation used for tuning.

The results compare forecasting procedures matched on those declared
dimensions. Equal \(W\) alone does not isolate architecture from all other
information and tuning choices. This qualification does not assert future-label
leakage or replace any result. The broader treatment of estimation and data
use as components of a forecasting method is consistent with
\href{https://doi.org/10.1111/j.1468-0262.2006.00718.x}{Giacomini and White
(2006)}; the present information-footprint distinction concerns the specified
procedures in this study.

\subsection*{D. Sorting crossing horizon reports weakly improves aggregate pinball loss}

At one issue origin, let the realized short- and long-horizon peaks satisfy
\(y_s\le y_l\). Suppose the already available reports cross: the short-horizon
report is \(a\), the long-horizon report is \(b\), and \(a>b\). Use the same
quantile level \(0<\tau<1\), common units and the equally weighted sum of the
two pinball losses. Sorting the reports into horizon order replaces
\((a,b)\) by \((b,a)\).

The identity \(\rho_\tau(y-c)=\tau(y-c)+(c-y)_+\) gives the original sum
minus the sorted sum as
\[
\begin{aligned}
D={}&\rho_\tau(y_s-a)+\rho_\tau(y_l-b)
      -\rho_\tau(y_s-b)-\rho_\tau(y_l-a)\\
 ={}&(a-y_s)_+-(b-y_s)_+-(a-y_l)_++(b-y_l)_+\\
 ={}&\int_b^a\bigl[\mathbf1\{u>y_s\}-\mathbf1\{u>y_l\}\bigr]\,du\\
 ={}&\operatorname{length}\bigl([b,a]\cap[y_s,y_l]\bigr)\ge0.
\end{aligned}
\]
The inequality is strict exactly when the intersection has positive length.
The transformation uses only the issued reports; future outcomes enter the
proof of its score property, not the transformation. If both reports satisfy
the common observed floor \(B\), their sorted versions still satisfy it.
The same statement holds after averaging the sum over common issue origins.

\textbf{Correction to the earlier Supplement.} Its statement that raising a
long-horizon report can worsen that report's pinball loss is correct. Its
following assertion that target ordering does not identify a universally
beneficial repair is too broad when read as an assertion about aggregate
equal-weight pinball loss: the identity above provides such a weak improvement.
The earlier counterexample rules out a raise-only marginal-loss guarantee,
rather than this sorting result. The protected Supplement and its crossing
counts are retained; this note explicitly supplies the qualification.

This aggregate result provides no guarantee for each horizon's loss
separately, arbitrary unequal weights, differently normalized horizon scores,
breach counts or conditional coverage. No sorted empirical forecasts or new
financial loss summaries are reported here. Related rearrangement work gives
useful context: \href{https://arxiv.org/pdf/0704.3649}{Chernozhukov,
Fern\'andez-Val and Galichon, Proposition 4} establishes an \(L^p\)
estimation-error improvement for a quantile curve rearranged over its
probability index. That proposition is not the fixed-\(\tau\), cross-horizon,
samplewise pinball identity proved here. No new general rearrangement theorem
is claimed.

\subsection*{E. Horizon ordering uses a common lower-inverse selection}

For nested peaks under the same issue-information law,
\(Y_{t,H_s}\le Y_{t,H_l}\) implies
\(F_s(z\mid\mathcal F_t)\ge F_l(z\mid\mathcal F_t)\) for every \(z\).
At the same level \(\tau\), using the common lower-inverse convention
\[
q^-_\tau(H\mid\mathcal F_t)
=\inf\{z:F_H(z\mid\mathcal F_t)\ge\tau\}
\]
therefore yields
\(q^-_\tau(H_s\mid\mathcal F_t)\le
q^-_\tau(H_l\mid\mathcal F_t)\). Atoms and flat CDF intervals are allowed.
This is the canonical selection already used in the Supplement.

The full quantile functional can instead be set-valued. For example, if the
two peaks coincide with a variable taking zero with probability \(\tau\) and
one with probability \(1-\tau\), every report in \([0,1]\) minimizes expected
pinball loss. Choosing different members of this interval at different
horizons can produce crossing optimal reports. Consequently crossing is a
coherence diagnostic relative to a declared selection convention; by itself
it does not establish marginal misspecification.
\href{https://arxiv.org/pdf/1503.08195v2}{Ehm et al., Sections 2.1--2.2}
define the potentially set-valued quantile functional and its consistent
scoring functions. This selection qualification is distinct from both the
aggregate sorting identity and the strict-tail coverage inequality for a true
conditional lower-inverse quantile.


\subsection*{F. Observation and calibration contracts}

In the weekly experiment, the current-position floor is observed under a
stipulated economic-date replay of one downloaded vintage. This convention
makes the current and prior snapshot values available at each issue index.
It is distinct from the historical information of a public observer: the
\href{https://www.newyorkfed.org/markets/counterparties/primary-dealers-statistics}{New
York Fed source page} describes publication in the following week, and
later versions may revise earlier observations. Delaying only the predictors
would not repair a public-information interpretation while retaining an
unreleased current-position floor. No release-vintage reconstruction or
changed financial experiment is supplied by this clarification.

The main delayed-calibration comparison concerns the score advantage of the
wrapped regression over the wrapped historical reference. Each wrapper's
saved mean loss is higher than that of its own unwrapped base. Thus a retained
relative advantage and closer observed breach frequencies coexist with a
score cost relative to the bases; this is not a claim of joint absolute
improvement. The validation warmup and calibration-parameter choice are the
already documented retrospective second-stage selection, not a fresh
historical issue-by-issue model-selection exercise.

The earlier Supplement's sentence about state carrying across the unissued
gap refers to the selected ACI candidate's calibration state. It should not
be read as a claim that the headroom-SAV fitted state is carried from validation
into evaluation. That recursion is refitted on its evaluation segment's
specified completed labels and refit grid. These distinctions follow the
frozen implementations and existing saved outputs; no canonical ACI theorem,
CAViaR-family result, new score or statistical comparison is imported.
